\scanpage{38}
\mspage{34}
% archive sheet number: 37
% upper margin note, enclosed by a long curved bracket
\margin{Terminologie : atelier, magasin, strate, multistrates.}
% source lines 1--7
\uline{Digression 4}\quad Ainsi, de plus en plus, l'accent principal dans le travail de fondements, se porte sur $M$, plutôt que sur $\unsure{\mathcal F|\mathscr F|\mathfrak F}$. Du coup, cela suggère aussi un changement de terminologie. Le \deleted{terme}\inserted{nom} \unsure{initial} principal, d'\uline{atelier}, \uline{devrait être donné à} $M$.
% source lines 7--18
Cet atelier travaille avec les objets qui figurent dans $M$, les \inserted{strates \unsure{ou}} multistrates, pour constituer avec eux des \uline{figures}. L'ensemble de ces figures peut-être considéré comme constituant un \uline{magasin attenant à l'atelier} (plutôt que l'inverse) --- plusieurs magasins peuvent être attenants au même atelier, comme on vient de le rappeler.
% source lines 19--23
Dans le cas d'un magasin \unsure{compactologisé}, on pourrait appeler peut-être $M_c$, formé des (multi)strates compactes, le \uline{coeur} de l'atelier.
